\subsection{All abelian preimages, their horizontal obstruction, and the filled elliptic target}

\begin{statement}{E.87}{Proved}{The rank-seven relative lattice and an evaluated free wild obstruction}
Retain E.85's actual source marking, its section
$(u_0,d)$, and the intrinsic abelian class $\alpha_0$.
The actual map from $V=N^{\rm ab}$ to the joint
abelian lattice $\mathcal J_1$, in its coordinates
$(L,p,m)$, is
\[
 R=\begin{pmatrix}
 1&-1&0&0&0&-1&-1&-1&0&0\\
 2&-1&0&0&0&0&0&0&0&-3\\
 1&-1&0&0&0&0&0&0&-1&-1
 \end{pmatrix}.
 \tag{E87.1}
\]
Its actual local cohomology map is onto, with the
following complete kernel and fibre type:
\[
 \begin{gathered}
 H^1(\mathbb Q_3,V)=\mathbb Z_3^{15}\oplus C_3^3
   \twoheadrightarrow
 H^1(\mathbb Q_3,\mathcal J_1)=\mathbb Z_3^5\oplus C_3,\\
 \ker H^1(R)\simeq\mathbb Z_3^{10}\oplus C_3^3.
 \end{gathered}
 \tag{E87.2}
\]
All three source torsion coordinates map to zero;
the target torsion coordinate is reached from a
free coordinate. In particular this exact sequence
has no $\mathbb Z_3$-linear section. The following
actual inverse coordinates include its finite carry.

Use the retained Smith matrix $L_1$ for
$(I-T_1)(I+J_1)$, with diagonal
$(1,3,3,3,0^6)$. Write a horizontal class as
\[
 L_1U=(0,\epsilon_1,\epsilon_2,\epsilon_3,q_1,\ldots,q_6),
                     \qquad\epsilon_i\in C_3.
 \tag{E87.3}
\]
Its target horizontal coordinates are
\[
 L_u=-q_2+q_3-q_4-q_5,\quad
 m_u=q_3-q_5-q_6,\quad
 \epsilon=2q_3-q_5\pmod3.
 \tag{E87.4}
\]
For a prescribed $(L_u,m_u,\epsilon)$, take
$\epsilon\in\{0,1,2\}$, retain
$q_1,q_3,q_4,k\in\mathbb Z_3$ and all three
source torsion coordinates, and set
\[
 q_5=2q_3-\epsilon-3k,\qquad
 q_6=-q_3+\epsilon+3k-m_u,\qquad
 q_2=-q_3-q_4+\epsilon+3k-L_u.
 \tag{E87.5}
\]
The inverse recovers $k=(2q_3-q_5-\epsilon)/3$.
For the relative coordinate use E.75's basis
\[
 A=\theta(3f_1+f_2)+a_0A_0+a_1A_1+a_2A_2
             +b_0B_0+b_1B_1+b_2B_2+c_1C_1+c_2C_2.
 \tag{E87.6}
\]
Its image is
\[
  (L_a,p_a,m_a)=(2\theta-b_0-b_1-b_2,
                  5\theta-3c_2,2\theta-c_1-c_2).
 \tag{E87.7}
\]
Retain freely $a_0,a_1,a_2,b_1,b_2,c_2$ and solve
\[
 \theta=(p_a+3c_2)/5,\quad
 c_1=2\theta-c_2-m_a,\quad
 b_0=2\theta-b_1-b_2-L_a.
 \tag{E87.8}
\]
Equations (E87.3)--(E87.8), with the literal inverse
Smith matrix, are mutually inverse coordinates on
every fibre in (E87.2).

There is also an exact coefficient-kernel description.
Put $K_1=\ker R$ and use its following seven basis vectors:
\[
 \begin{aligned}
 k_1&=f_1+2f_2-B_0-C_1,&
 k_2&=-3f_2+3B_0+2C_1+C_2,\\
 k_3&=A_0,&k_4&=A_1,&k_5&=A_2,\\
 k_6&=B_1-B_0,&k_7&=B_2-B_0.
 \end{aligned}
 \tag{E87.9}
\]
If $K$ is this ten-by-seven matrix, its left inverse
selects rows $f_1,C_2,A_0,A_1,A_2,B_1,B_2$.
The exact operators are $T_K=K^\ell T_1K$ and
$J_K=K^\ell J_1K$, saved as integral matrices, and
\[
 \begin{gathered}
 H^1(\mathbb Q_3,K_1)
    \xrightarrow{\ \sim\ }\ker H^1(R),\\
 \mathcal P_K=K_1/(I-T_K)K_1^{J_K}
       \simeq\mathbb Z_3^4\oplus C_3^3,\\
 \mathcal K_K=
  \ker((I-J_K)(T_K-7I))
       =\mathbb Z_3(9k_1+5k_2)
                    \oplus\bigoplus_{i=3}^7\mathbb Z_3k_i,\\
 H^1(\mathbb Q_3,K_1)=\mathcal P_K\oplus\mathcal K_K.
 \end{gathered}
 \tag{E87.10}
\]
These are the coefficient and cohomology maps of
the actual exact sequence $0\to K_1\to V\to\mathcal J_1\to0$.

\emph{A complete larger intrinsic locus.}
Consider all intrinsic abelian classes whose relative
value is still $\bar d=e_{B_0}$ and whose joint
abelian image is $\beta_0$. They are exactly
$\alpha_0+\mathcal P_K$. For
$v=\sum_{i=1}^7x_i k_i$, their full class-two
obstruction is the evaluated homomorphism
\[
 \begin{gathered}
 \delta_{B_0}:\mathcal P_K\longrightarrow
    \mathcal B_2=\mathbb Z_3\oplus C_3^{11}\oplus C_9^4,
        \qquad [v]\longmapsto[7[v,e_{B_0}]],\\
 \operatorname{im}\delta_{B_0}\simeq C_3^3,
       \qquad\ker\delta_{B_0}\simeq\mathbb Z_3^4.
 \end{gathered}
 \tag{E87.11}
\]
There are exactly three scalar conditions, all modulo three:
\[
 x_1-2x_6-x_7=0,\qquad x_2=0,\qquad
 x_4-x_5-x_6-2x_7=0.
 \tag{E87.12}
\]
Their full integral solution lattice is the column
lattice of
\[
 Q_K=\begin{pmatrix}
 3&0&0&0&0&2&1\\0&3&0&0&0&0&0\\
 0&0&1&0&0&0&0\\0&0&0&3&1&1&2\\
 0&0&0&0&1&0&0\\0&0&0&0&0&1&0\\
 0&0&0&0&0&0&1
 \end{pmatrix},\qquad\det Q_K=27.
 \tag{E87.13}
\]
The exact Smith matrix of
$Q_K^{-1}(I-T_K)(I+J_K)$ has three unit factors
and four zero factors. It supplies both inverse
coordinates on the free four-dimensional zero locus.

All these abelian classes have actual intrinsic
class-two lifts. Their entire class-two section
locus maps to the same index-nine part
\[
                  \mathbb Z_3\oplus3\mathbb Z_3\oplus0
                      \subset\mathfrak G_{\beta_0}
 \tag{E87.14}
\]
computed in E.85. Every point of this image has,
within this whole inertia-fixed locus, a complete
geometric-conjugacy fibre
\[
                     \mathbb Z_3^{48}\oplus C_3^9\oplus C_9^5.
 \tag{E87.15}
\]
The four additional coordinates are the actual
kernel in (E87.11); the other forty-four and the
fourteen cyclic coordinates are E.85's full central
fibre coordinates. The proof specifies the particular
central corrections, the joint normalization and the
remaining actual geometric conjugators.

\emph{An evaluated quotient for all intrinsic abelian inputs.}
Fill the eight cusps other than $O=C_0$ and put
\[
 F_E=\pi_1^{(3)}(E\setminus\{O\}),\qquad
 C_E=F_E/\gamma_3F_E.
 \tag{E87.16}
\]
There is the actual equivariant quotient
\[
 N_2\twoheadrightarrow C_E,\qquad
 \ker=\frac{\langle\!\langle d_i:i\ne O\rangle\!\rangle
                         \gamma_3N}{\gamma_3N}.
 \tag{E87.17}
\]
Its abelian map sends $f_1,f_2$ to the two elliptic
handle generators and sends the eight other basis
vectors to zero. Its central map is projection to
the coefficient of $[f_2,f_1]$, precisely E.74's
free $\mathcal B_2$ coordinate.

The actual arithmetic action on $C_E$ is completely
evaluated after the following explicit central
coordinate morphism. Choose handle generators
$g_1,g_2$ with the stated abelian images, and put
$X_i=\log g_i$, $Z=[X_2,X_1]$. The rational
involution $\iota[x:y:z]=[y:x:z]$, with its
coherent tangential comparison below, has
\[
 \iota X_i=-X_i+b_iZ,\qquad \iota Z=Z.
 \tag{E87.18}
\]
The coefficients $b_i$ here are the actual central
coordinates of $\iota(g_i)g_i$. Retain them in the
literal invertible change of basis
\[
 X_i^\sharp=X_i-\tfrac12b_iZ,\qquad
 g_i^\sharp=g_i[g_2,g_1]^{-b_i/2}.
 \tag{E87.19}
\]
For old group Hall coordinates
$g_1^xg_2^y[g_2,g_1]^n$, the new Lie central
coordinate and its inverse are
\[
 r^\sharp=n-xy/2+(b_1x+b_2y)/2,
 \qquad n=r^\sharp+xy/2-(b_1x+b_2y)/2.
 \tag{E87.20}
\]
Thus the comparison preserves the E.37 abelian
marking and retains every central correction.

In coordinates $(X_1^\sharp,X_2^\sharp,Z)$ the
genuine arithmetic operators are exactly
\[
 T=(A_E,7),\quad J=(C_E^{\rm mat},-1),\qquad
 A_E=\begin{pmatrix}1&-3\\3&-2\end{pmatrix},\quad
 C_E^{\rm mat}=\begin{pmatrix}1&-1\\0&-1\end{pmatrix}.
 \tag{E87.21}
\]
The first-degree section coordinates of this
quotient are $(U_1,t,\epsilon)\in\mathbb Z_3^2\times C_3$,
where its relative value is $(3t,t)$ and its
horizontal value is $(U_1,U_2)$ with
$\epsilon=U_2\bmod3$. Its complete class-two
lifting image is
\[
 \begin{gathered}
              -7t(U_1+21t/2)=0,\\
              t=0\quad\hbox{or}\quad U_1=-21t/2.
 \end{gathered}
 \tag{E87.22}
\]
Each point of this exact closed locus has fibre
$\mathbb Z_3^2$. Its inverse takes arbitrary
central Lie coordinates $(r_u,r_a)$ and uses
\[
      u=(U_1,\epsilon,r_u),\qquad
      a=(3t,t,r_a)
 \tag{E87.23}
\]
as E.77's two values in $C_E$.
For an unnormalized representative the actual fixed
conjugator is $\exp(bX_1^\sharp)$, where
$b=(U_2-\epsilon)/3$. Its full central corrections are
\[
 r_u\longmapsto r_u+\frac{3b^2+3bU_1-2bU_2}{2},
 \qquad r_a\longmapsto r_a-bt.
 \tag{E87.24}
\]
There is no remaining lower stabilizer. These
coordinates therefore include the complete
geometric-conjugacy fibres.

For any intrinsic abelian class, projection through
(E87.17) makes (E87.22) its exact free
class-two obstruction. Consequently every
intrinsic class-two section lies over these two
evaluated elliptic branches. The fifteen finite
coordinates of $\mathcal B_2$ for general relative
inputs remain separate from this evaluated free
coordinate; on the horizontal locus they have
already been evaluated by (E87.11)--(E87.13).

\emph{Manuscript: \S51.16, folios733v--736v,
scans172--178, formulas(371)--(379), Part E(c),(d)
and follow-up(4). All intrinsic, joint and filled
maps in this result are the actual displayed
quotient or coefficient maps.}
\end{statement}

\begin{proof}
\emph{The complete first-degree comparison.}
The rows in (E87.1) are obtained from the four
actual unit characters in E.83 as
$X_h-Y_h$, $X_h$ and $(X_h+X_{\bar h})/3$.
They are integral, and their columns $f_2,B_0,C_1$
have determinant a unit. Hence $R$ is onto.
Direct multiplication gives
\[
          RT_1=7R,\qquad RJ_1=J_{\mathcal J_1}R.
 \tag{E87.25}
\]
The actual E.77 source, reduced to any abelian
lattice $M$, has coordinates
\[
 \frac{M}{(I-T)M^J}\oplus
          \ker((I-J)(T-7):M\to M),
 \tag{E87.26}
\]
with the source values $(U,A)$ and the complete
fixed geometric gauge. This follows either
directly from its relator or from E.75's exact
cochain complex, using the specified source
marking throughout.

The first-degree Smith matrix has the stated
diagonal and the following literal coefficient
product:
\[
 RL_1^{-1}=\begin{pmatrix}
 0&0&0&0&0&-1&1&-1&-1&0\\
 -6&-6&6&6&0&0&2&0&-1&-3\\
 0&0&0&0&0&0&1&0&-1&-1
 \end{pmatrix}.
 \tag{E87.27}
\]
The target fixed lattice is generated by its
$p$ coordinate, and its gauge is multiplication
by six on that coordinate. Thus (E87.27) gives
exactly (E87.4). It proves that all source torsion
maps to zero. Equations (E87.5) give all
horizontal preimages and their inverse, including
the residue carry. The actual relative kernel
on $V$ is generated by $3f_1+f_2$ and the eight
cusp vectors. Multiplication by $R$ gives
(E87.7), and (E87.8) proves the full relative
surjectivity with its inverse. This proves
(E87.2) and all the claimed fibre coordinates.
A linear section would send the target element
of order three to a source element of order
three, contradicting the computed zero map
on source torsion. This proves the stated
non-splitting assertion.

Solving the three rows of $R$ gives exactly the
seven vectors (E87.9) and the displayed left
inverse. They provide a saturated coefficient
kernel. In these coordinates its operators are
\[
 T_K=\begin{pmatrix}
 -5&9&0&0&0&0&0\\-3&4&0&0&0&0&0\\
 5&-9&7&0&0&0&0\\2&-12&0&7&0&0&0\\
 -4&3&0&0&7&0&0\\1&-6&0&0&0&7&0\\
 -8&6&0&0&0&0&7
 \end{pmatrix},\qquad
 J_K=\begin{pmatrix}
 -1&3&0&0&0&0&0\\0&1&0&0&0&0&0\\
 0&0&-1&0&0&0&0\\1&0&0&0&-1&0&0\\
 -1&3&0&-1&0&0&0\\1&0&0&0&0&0&-1\\
 -1&3&0&0&0&-1&0
 \end{pmatrix}.
 \tag{E87.28}
\]
The horizontal Smith factors are $3,3,3,0^4$.
The relative equation is the single saturated
condition $5x_1=9x_2$, giving the relative basis
in (E87.10). Its inverse divides $x_2$ by the
unit five. Since $H^0(\mathbb Q_3,\mathcal J_1)=0$
and $H^1(R)$ is onto, the exact cohomology sequence
identifies $H^1(K_1)$ with the full kernel.
This also proves that the displayed coefficient
description includes every abelian preimage.

\emph{The complete horizontal obstruction.}
Fix the actual section $(u_0,d)$ of E.85 and
choose a continuous geometric Hall lift $g_v$
of $v\in K_1$ into $N_2$. In class two the
commutator depends only on abelian coordinates.
The full cusp norm identity therefore gives
\[
 (u_0g_v)T(d)(u_0g_v)^{-1}d^{-7}
                         =7[v,e_{B_0}]\quad\hbox{in }A_2.
 \tag{E87.29}
\]
Replacing the relative value by $d w_a$ adds
$(T_2-7I)a$. Hence the exact source condition is
\[
 (I-J_2)\bigl(7[v,e_{B_0}]+(T_2-7I)a\bigr)=0.
 \tag{E87.30}
\]
E.77's factor of one half identifies its
obstruction with the class in (E87.11).
For a horizontal gauge vector $b\in K_1^{J_K}$,
the bracket identity
\[
 7[(I-T_1)b,e_{B_0}]
            =-(T_2-7I)[b,e_{B_0}]
 \tag{E87.31}
\]
proves that this map descends through the full
abelian gauge. Its free coefficient is zero,
because every bracket in (E87.29) contains a
cusp vector.

The certificate computes its fifteen retained
E.74 coordinates and their moduli. All its
values have order dividing three, including
the components landing in $3C_9$. Dividing
those four coordinate subgroups by three,
the three selected target rows numbered
$4,8,15$ have an invertible three-by-three
matrix over $\mathbb F_3$. The entire matrix
factors through the three characters in
(E87.12), with that inverse. This proves
its exact image $C_3^3$ and exact kernel
lattice (E87.13). The actual transformed
gauge matrix has three unit factors, proving
that the cohomological zero locus is free
of rank four. Both inverse Smith maps are
retained, rather than just these group orders.
Explicitly, if
$L_Q(Q_K^{-1}(I-T_K)(I+J_K))R_Q$
has diagonal $(1,1,1,0,0,0,0)$,
the chosen representative of
$(q_1,q_2,q_3,q_4)\in\mathbb Z_3^4$ is
$v=KQ_KL_Q^{-1}(0,0,0,q_1,q_2,q_3,q_4)^{\mathsf t}$.

Here is the actual particular lift for every
point of that lattice. Use the exact E.75
Smith system for
\[
 D^-=(I-J_2)(T_2-7I):A_2\longrightarrow A_2^-.
 \tag{E87.32}
\]
Let $B^-$ be its minus-eigenspace basis,
$(B^-)^\ell$ its left inverse, and
$L^-D^-_{\rm coord}R^-=\Delta^-$ its exact
Smith system. The first twenty-three
diagonal entries $d_i$ are nonzero; the
remaining row is zero. For $v=KQ_Kx$ put
\[
 \begin{gathered}
 \omega=(I-J_2)7[v,e_{B_0}],\qquad
 y=L^-(B^-)^\ell\omega,\\
 a(v)=-R^-_{*,1:23}
           \operatorname{diag}(d_i^{-1})y_{1:23}.
 \end{gathered}
 \tag{E87.33}
\]
Every entry of this forty-five-by-seven
matrix is in $\mathbb Z_{(3)}$; its exact
product with $D^-$ is $-\omega$. These
are verified rational matrix identities.
Thus (E87.33) is an actual integral lift,
not an inverse-limit inference from finite
vanishings.

The image $F_2a(v)$ is diagonal, because
$F_2[v,e_{B_0}]=0$ and (E87.30) holds.
It lies in $L_2$, so write it as $(3s,3s)$.
Replace $a(v)$ by $a(v)-s a_*$, with the
literal vector of E.85. This preserves
(E87.30) and makes its joint central image
zero. The resulting exact correction matrix
and all its columns are saved.

It remains to normalize the horizontal joint
value without changing its intrinsic abelian
coordinate. The pair $g_v$ has zero joint
abelian image, hence has central image
$(c^n,c^m)$. The section just constructed has
relative joint value $(y,y)$. Its joint
horizontal class is E.80's family with
parameters $(2+n,3+m)$. Applying E.83 to
this actual intrinsic section proves
$n+m\in3\mathbb Z_3$. The literal Hall
right inverse (E85.22) therefore gives
$z_v$ with $F_2z_v=(n,m)$. Replacing $g_v$
by $g_vw_{-z_v}$ makes its whole class-two
joint image one; it leaves its abelian
coordinate and (E87.29) unchanged.
This is an actual continuous normalization:
the joint coordinates of the chosen Hall
lift are continuous, and division of
$n+m$ by three is taken on its proved
image. No value of a geometric path is
replaced by an assumed zero.

The two resulting values now give an
intrinsic section mapping exactly to the
E.85 origin. E.85's central splitting and
kernel coordinates supply every target
point in (E87.14) and every central fibre.
Conversely any section in this inertia-fixed
abelian locus has zero obstruction (E87.11),
so its horizontal abelian class belongs to
the computed free four-dimensional kernel.
The abelian stabilizer is zero, as in E.85.
Its unique aligning abelian conjugator is
lifted to the actual fixed subgroup; the
gauge uses its literal $g u T(g)^{-1}$ and
$gag^{-1}$, retaining any nonlinear central
correction. The remaining gauge is precisely
the central one of E.85. These observations
prove both directions of the full fibre
description (E87.15).
For a completely specified fixed lift of
an aligning vector $b$, choose its ordered
Hall lift $g_b$ and put
$\delta=J(g_b)g_b^{-1}\in A_2$.
Then $J\delta=-\delta$ and
$g=\delta^{1/2}g_b$ is fixed. Thus this
alignment too is an actual continuous
coordinate map, with the indicated
central correction retained.

\emph{The rational involution and its exact tangent.}
The map $\iota[x:y:z]=[y:x:z]$ is a rational
order-two automorphism of the cubic fixing
$O=[1:-1:0]$. It interchanges $A_i$ with
$B_i$ and sends $C_i$ to $C_{-i}$, so it
preserves the eight-cusp filling kernel.
Here is an exact verification that it
is $[-1]$ on the elliptic curve with
origin $O$. Put $r=(x+y)/z$ and
$v=(x-y)/z$. The cubic equation gives
$v^2=-(r^3+4)/(3r)$, so the quotient
by this involution is the rational
$r$-line and has degree two. On
$\operatorname{Pic}^0(E)$, pullback
after the divisor norm is $1+\iota$.
It factors through
$\operatorname{Pic}^0(\mathbb P^1)=0$,
hence $1+\iota=0$. The Abel map
$P\mapsto[P-O]$ intertwines the
involution because $O$ is fixed.
This proves the asserted multiplication
map and its action $-I$ on the actual
elliptic Tate module.

The based comparison is fully specified.
Near $O$ write $s=z/x$ and
$y/x=-(1+s^3)^{1/3}$ with constant term
$-1$. Then
\[
 \iota(s)=-s(1+s^3)^{-1/3},\qquad
 s^{1/3^n}\longmapsto
     -s^{1/3^n}(1+s^3)^{-1/3^{n+1}}.
 \tag{E87.34}
\]
All unit roots in this formula have constant
term one and rational coefficients. The
constant roots of $-1$ are all $-1$, and
the displayed maps are compatible under
cubing. Moreover
$\iota(s)^3=-s^3/(1+s^3)$, so
$1+\iota(s)^3=(1+s^3)^{-1}$.
Applying the displayed root map twice
therefore gives exactly $s^{1/3^n}$.
Thus the comparison is an actual based
involution and commutes with the entire
arithmetic action. It descends through the
filling map, which is defined over
$\mathbb Q$ with the same tangent.

For the filling map the complex
fundamental-polygon presentation has
two handle generators and nine
peripheral generators, with one
relation equating their product
to the handle commutator. Killing
the eight specified peripherals
leaves the free group on the two
handles. Riemann existence and
its pointed comparison
\cite[Expos\'e XII, Th\'eor\`eme~5.1 and Corollaire~5.2]{E_SGA1}, followed by the
universal property of pro-three
completion, therefore give exactly
the stated closed normal kernel
and the surjection on geometric
groups. The images of the two chosen
intrinsic handle generators form its
abelian basis; they form an actual free
basis by the pro-three Frattini criterion
and the Hopfian property. This proves
(E87.17), including its abelian and
central maps. The quotient kills all
the other second-degree Hall vectors.

\emph{The actual central coordinate morphism.}
For a class-two pro-three group, the
integral Baker--Campbell--Hausdorff law
uses only division by two. In the stated
rank-two basis it is the explicit law
\[
 (X,Y,r)(X',Y',r')
   =(X+X',Y+Y',r+r'-(XY'-YX')/2).
 \tag{E87.35}
\]
It follows directly from the class-two
Hall collection law, so logarithm and
exponential here are mutually inverse
continuous integral coordinate maps.
Since $\iota$ is $-I$ on the elliptic
Tate module and fixes its central wedge,
its general action has the form (E87.18).
Applying (E87.19) makes it exactly
$(-I,+I)$. Direct substitution gives
the inverse formula (E87.20).

Every genuine arithmetic automorphism
commutes with this involution. On the
new handle basis it has the form
\[
 X_i^\sharp\longmapsto
       \sum_j A_{ji}(g)X_j^\sharp+\eta_i(g)Z.
 \tag{E87.36}
\]
Commuting with $(-I,+I)$ gives
$2\eta_i(g)=0$, hence $\eta_i(g)=0$
in the torsion-free central lattice.
Its action on $Z$ is $\det A(g)=\chi(g)$.
This proves the exact graded action
through the displayed coordinate
morphism, with no replacement of the
original based group. In the old basis
the same computation gives the actual
cross row
\[
             \eta(g)=\bigl(\chi(g)b-bA(g)\bigr)/2,
                   \qquad b=(b_1,b_2).
 \tag{E87.37}
\]
Thus the old and new arithmetic actions
are related by the full central formula,
not just by their abelian matrices.
The actual E.37 matrices now give (E87.21).

\emph{The complete filled local section set.}
The abelian E.77 relation is
$(I-C_E^{\rm mat})(A_E-7I)A=0$.
Its exact kernel is $A=(3t,t)$.
The fixed abelian gauge is generated
by $(1,0)$, and $(I-A_E)(b,0)=(0,-3b)$.
This gives the three first-degree
coordinates used in the statement.

Let $u=(U_1,U_2,r_u)$ and
$a=(3t,t,r_a)$ in (E87.35).
Direct multiplication of the actual
source norm word gives
\[
 uT(a)u^{-1}a^{-7}
       =(-21t,0,-7t(U_1+21t/2)).
 \tag{E87.38}
\]
Its abelian coordinate is fixed by
$C_E^{\rm mat}$, and its central
coordinate is negated by $J$.
The complete source relation is
therefore exactly (E87.22). The
central obstruction normalization
of E.77 divides the doubled central
defect by two, giving the displayed
free obstruction with its sign.
Since $\mathbb Z_3$ is an integral
domain, its vanishing locus is
exactly the union of the two
stated branches. Both central
coordinates are unrestricted.

The fixed geometric conjugator
has coordinates $(b,0,0)$: its
central coordinate must vanish
because $J$ negates it. Applying
the full group law to
$g u T(g)^{-1}$ and $gag^{-1}$
gives respectively
\[
 \begin{aligned}
 &(U_1,U_2-3b,
       r_u+(3b^2+3bU_1-2bU_2)/2),\\
 &(3t,t,r_a-bt).
 \end{aligned}
 \tag{E87.39}
\]
Taking $b=(U_2-\epsilon)/3$
proves (E87.24) and its inverse.
There is no nonzero abelian fixed
vector because $\det(I-A_E)=9$;
there is no central fixed vector
because $7-1$ is injective.
Thus no further stabilizer acts
on the two central coordinates.
This proves the full fibre and
homeomorphism assertions.

Finally the central projection
in (E87.17) is the exact free
coordinate used in E.74. The
natural obstruction square for
these actual central extensions,
proved by projecting a continuous
Hall lift and its defect as in
E.83, identifies (E87.22) with
the free component of every
intrinsic obstruction. This
justifies each intrinsic exclusion
through the actual filled map.

All matrix divisions above take
place on proved integral images.
The retained finite Smith systems
contain both inverse maps, and
the displayed particular lift
is an exact rational identity
over $\mathbb Z_{(3)}$. All
coordinate spaces are compact;
their explicit continuous inverses
give the asserted compact
topologies. The actual code checks
the full first-degree maps, all
$2,187$ residue inputs for the
horizontal obstruction, every
particular correction column,
the coherent formal-disc square,
the central averaging morphism,
the symbolic filled norm word
and $729$ complete gauge/inverse
identities. The source and input
hashes and all actual outputs
are retained.
\end{proof}
